\sect{14. Finding your way and hints}
\textbf{Your mission, step by step.} Explore each floor, find routes deeper into the dungeon, collect the required keys and pieces, and keep the party able to fight and heal. The final goal is to find and defeat Xanathar. Secret rooms and optional quests can reward exploration, but finding every secret is different from completing the main mission.

\textbf{Automap.} The adventure appears on the V9990 output and the map on the MSX output. Connect two displays or switch monitor inputs. In openMSX, choose MSX as the video source for the map and GFX9000 to return to the adventure.

North is at the top. The white arrow shows where the party stands and which way it faces. Black squares are unexplored. The map marks visible doors in orange, stairs up in cyan and stairs down in yellow. Return to a place to see changes there. To reveal an illusion wall as a passage, walk through it.

The game saves the map with your progress. Open \texttt{CAMP} to pause the adventure while consulting it.

\textbf{Browsing floors.} \key{F2} shows the previous floor, \key{F3} the next and \key{F4} returns to your current floor. This also works in camp. Visible open pits are red and portals magenta. Keep separate notes about inscriptions and portal destinations.

\subsection*{Your first few minutes}
\begin{enumerate}
\item Put the best-protected characters in front. Check their weapons are in their hands, not their backpacks.
\item Open inventories and locate the book, symbol, food and ammunition. Read the commission letter.
\item Select healing and some attack spells; rest somewhere safe to recover them.
\item Explore a short stretch, watch the compass and walls, collect useful items and try visible buttons.
\item Save before your first descent. Keep a route back to a familiar place.
\end{enumerate}

\subsection*{Hints, from gentle to more helpful}
\begin{itemize}
\item \textbf{If it looks like a dead end:} inspect side walls, read inscriptions and check recesses and small buttons. Not every passage opens with a key.
\item \textbf{If a mechanism will not respond:} consider the object, its shape and where to place it. Carrying a piece in your backpack is not the same as using it in the correct socket.
\item \textbf{If a door closes again:} look for a pressure plate and consider whether it needs weight. An ordinary object may be more useful on the floor than in a backpack.
\item \textbf{If you reappear somewhere familiar:} there may be a portal or a change in heading. Check the compass before counting more steps.
\item \textbf{If you find stone objects:} note their shapes and places with matching sockets. Think before using or abandoning them: they help connect routes.
\item \textbf{If you find bones or friendly characters:} listen before attacking and keep the remains. Some encounters provide companions, help or clues.
\item \textbf{If you lose a fight:} return with prepared spells, a clear retreat and healing in the right hands. Avoid being surrounded and withdraw before both front characters fall.
\end{itemize}

Do not jump into a hole without preparing for a fall. Leave supplies in a recognisable place and note where. Keep keys and unique pieces; discard duplicate ordinary weapons first if space runs out.

\clearpage
\sect{15. Frequently asked questions and quick reference}
\begin{reference}{Problem}{Useful check}
I cannot attack & Is the character conscious? Has the hand stopped being grey? Is an enemy ahead, and is the character in front for a melee blow?\\
The bow will not fire & Put arrows in the other hand and spares in the quiver. Check that the character can use the bow.\\
I cannot equip an item & Check its slot, class restrictions and cursed equipment that prevents a change. An occupied hand also restricts two-handed weapons.\\
No spells appear & Check the book or symbol, character, spell level, CAMP selection and completed rest. Learned does not mean prepared.\\
Magic needs a free hand & Put away an item from the other hand. Some spells create an attack you then use from that hand.\\
Healing does not cure poison & These are different effects: antidote or Neutralize Poison first, HP healing afterwards.\\
I cannot rest & Look for nearby enemies, remove the danger and check food and party conditions.\\
I cannot drop a member & You need at least five characters and cannot choose the first.\\
A magic item behaves unexpectedly & Check its effect and how to use it in sections 11--13.\\
I cannot find my save & Check the selected language. Each language has its own save. Save your game before choosing \texttt{Quit Game}.\\
The adventure is silent & The opening fanfare uses the MSX's sound; adventure effects and introduction music require MoonSound. Check Sounds in Preferences.\\
\end{reference}
